\documentclass[
  spanish,
  aspectratio=169,
  xcolor={dvipsnames,table}
]{beamer}

\geometry{paperwidth=19.2cm,paperheight=10.8cm}

\usepackage{fontspec}
\setmainfont{Arial}
\setsansfont{Arial}
\usepackage[spanish,es-tabla,es-noshorthands]{babel}
\usepackage{array}
\usepackage{booktabs}
\usepackage{graphicx}
\usepackage{ragged2e}
\usepackage{tabularx}
\usepackage{amsmath,amssymb}
\usepackage{tikz}
\usepackage{hyperref}

\definecolor{uanlblue}{RGB}{0,55,108}
\definecolor{uanlgold}{RGB}{218,165,32}
\definecolor{tableblue}{RGB}{233,240,247}
\definecolor{darkcharcoal}{RGB}{33,37,41}

\setbeamercolor{palette primary}{bg=uanlblue,fg=white}
\setbeamercolor{palette secondary}{bg=uanlgold,fg=uanlblue}
\setbeamercolor{structure}{fg=uanlblue}
\setbeamercolor{title}{fg=white,bg=uanlblue}
\setbeamercolor{frametitle}{fg=uanlblue,bg=white}
\setbeamerfont{frametitle}{series=\bfseries,size=\large}

\setbeamertemplate{navigation symbols}{}
\setbeamertemplate{footline}{
  \leavevmode%
  \hbox{%
  \begin{beamercolorbox}[wd=.333333\paperwidth,ht=2.5ex,dp=1.1ex,left,leftskip=1cm]{palette primary}%
    \fontsize{7}{9}\selectfont UANL \textbar\ Facultad de Agronomía
  \end{beamercolorbox}%
  \begin{beamercolorbox}[wd=.333333\paperwidth,ht=2.5ex,dp=1.1ex,center]{palette primary}%
    \fontsize{7}{9}\selectfont Álgebra Lineal — PIA
  \end{beamercolorbox}%
  \begin{beamercolorbox}[wd=.333333\paperwidth,ht=2.5ex,dp=1.1ex,right,rightskip=1cm]{palette primary}%
    \fontsize{7}{9}\selectfont \insertframenumber{} / \inserttotalframenumber
  \end{beamercolorbox}}%
  \vskip0pt%
}

\begin{document}

\begin{frame}[plain]
  \begin{tikzpicture}[remember picture,overlay]
    \fill[uanlblue] (current page.south west) rectangle (current page.north east);
    \fill[uanlgold] ([yshift=1.2cm]current page.south west) rectangle ([yshift=1.45cm]current page.south east);
  \end{tikzpicture}
  \vspace{-0.5cm}
  \begin{center}
    \includegraphics[height=1.8cm]{UANL/assets-uanl/logos-institucionales.pdf}
    \vspace{0.4cm}
    
    {\fontsize{10}{12}\selectfont \textbf{\color{uanlgold}UNIVERSIDAD AUTÓNOMA DE NUEVO LEÓN \textbar\ FACULTAD DE AGRONOMÍA}}\\[0.2cm]
    {\fontsize{15}{18}\selectfont \textbf{\color{white}MODELACIÓN MATRICIAL DE FERTILIZACIÓN Y DINÁMICA POBLACIONAL CON MATRICES DE LESLIE}}\\[0.2cm]
    {\fontsize{11}{13}\selectfont \color{white}Producto Integrador de Aprendizaje (PIA) \textbar\ 25 Puntos}\\[0.5cm]

    \begin{tabular}{rl}
      \textbf{\color{uanlgold}Estudiante:} & \color{white}Carlos Mauricio Carvajal Coronado (Matrícula: 1977875) \\
      \textbf{\color{uanlgold}Docente Titular:} & \color{white}Miranda Vázquez José de Jesús \\
      \textbf{\color{uanlgold}Programa:} & \color{white}Ingeniero Agrónomo \textbar\ REF\#261013755137
    \end{tabular}
  \end{center}
\end{frame}

\begin{frame}{1. Justificación y Objetivos de la Modelación}
  \begin{columns}[T]
    \begin{column}{0.48\textwidth}
      \textbf{\color{uanlblue}Desafíos en la Citricultura Regional:}
      \begin{itemize}
        \item Alto costo de fertilizantes inorgánicos ($N, P, K$).
        \item Riesgo epidemiológico del HLB transmitido por \textit{Diaphorina citri}.
        \item Necesidad de modelos analíticos cuantitativos frente a recetas empíricas.
      \end{itemize}
    \end{column}
    \begin{column}{0.48\textwidth}
      \textbf{\color{uanlblue}Aportación del Álgebra Lineal:}
      \begin{itemize}
        \item \textbf{Sistemas Lineales ($A\mathbf{x} = \mathbf{b}$):} Solución exacta de dosis de mezcla física sin mermas ni excesos.
        \item \textbf{Matriz de Leslie ($L$):} Proyección estocástica de poblaciones por cohortes etarias.
        \item \textbf{Teorema Espectral:} Autovalor $\lambda_1$ para predecir colapso o explosión poblacional bajo control biológico.
      \end{itemize}
    \end{column}
  \end{columns}
\end{frame}

\begin{frame}{2. Formulación del Sistema de Fertilización ($A\mathbf{x}=\mathbf{b}$)}
  \begin{block}{Sistema Matricial de Insumo-Producto}
    $$\begin{pmatrix}
        0.46 & 0.18 & 0.13 \\
        0.00 & 0.46 & 0.00 \\
        0.00 & 0.00 & 0.44
    \end{pmatrix}
    \begin{pmatrix} x_1 \\ x_2 \\ x_3 \end{pmatrix}
    =
    \begin{pmatrix} 180 \\ 60 \\ 120 \end{pmatrix}$$
  \end{block}
  \begin{itemize}
    \item \textbf{Determinante:} $\det(A) = 0.46 \times 0.46 \times 0.44 = 0.093104 \neq 0$ $\implies$ Rango completo ($r=3$), solución única.
    \item \textbf{Dosis calculadas:}
      \begin{itemize}
        \item Nitrato de Potasio ($x_3$): $272.73\text{ kg/ha}$.
        \item Fosfato Diamónico 18-46-0 ($x_2$): $130.43\text{ kg/ha}$.
        \item Urea agrícola 46-0-0 ($x_1$): $263.19\text{ kg/ha}$.
      \end{itemize}
  \end{itemize}
\end{frame}

\begin{frame}{3. Dinámica Poblacional de Leslie: Formulación y Autovalores}
  \begin{columns}[T]
    \begin{column}{0.45\textwidth}
      \textbf{\color{uanlblue}Matriz de Leslie ($L$):}
      $$L = \begin{pmatrix}
        0 & 0 & 24 \\
        0.20 & 0 & 0 \\
        0 & 0.25 & 0.60
      \end{pmatrix}$$
      \vspace{0.2cm}
      \textbf{Ecuación característica:}
      $$\det(L - \lambda I) = \lambda^3 - 0.60\lambda^2 - 1.20 = 0$$
    \end{column}
    \begin{column}{0.51\textwidth}
      \textbf{\color{uanlblue}Comportamiento sin Control:}
      \begin{itemize}
        \item \textbf{Autovalor Dominante:} $\lambda_1 \approx 1.305 > 1.0$.
        \item \textbf{Tasa de crecimiento:} $30.5\%$ de incremento por ciclo de 10 días.
        \item \textbf{Estructura etaria asintótica ($\mathbf{w}_1$):}
          \begin{itemize}
            \item Huevos / ninfas tempranas: $82.8\%$
            \item Ninfas tardías: $12.7\%$
            \item Adultos reproductivos: $4.5\%$
          \end{itemize}
      \end{itemize}
    \end{column}
  \end{columns}
\end{frame}

\begin{frame}{4. Control Biológico con \textit{Tamarixia radiata} y Conclusiones}
  \begin{columns}[T]
    \begin{column}{0.48\textwidth}
      \textbf{\color{uanlblue}Simulación de Control Biológico:}
      \begin{itemize}
        \item Parasitismo selectivo reduce la supervivencia ninfal $s_2$ de $0.25$ a $0.05$.
        \item Nuevo polinomio característico:
          $$\lambda^3 - 0.60\lambda^2 - 0.24 = 0$$
        \item \textbf{Nuevo autovalor dominante:} $\lambda_1 = 0.898 < 1.0$.
        \item \textbf{Resultado:} La población decae a una tasa de $10.2\%$ por intervalo $\implies$ erradicación efectiva de la plaga.
      \end{itemize}
    \end{column}
    \begin{column}{0.48\textwidth}
      \textbf{\color{uanlblue}Conclusiones Clave:}
      \begin{enumerate}
        \item Se eliminó el desperdicio de fertilizantes mediante álgebra matricial.
        \item Se fundamentó científicamente el control biológico sin recurrir al abuso de insecticidas organofosforados.
        \item La matemática aplicada es una herramienta estratégica en la toma de decisiones agronómicas.
      \end{enumerate}
    \end{column}
  \end{columns}
\end{frame}

\end{document}

